\newpage
\begin{center}
    {\large \textbf{\tfgTitle}}

    \vspace{0.3cm}
    {\normalsize \tfgAuthors}

    \vspace{0.3cm}
\end{center}

\vspace{0.5cm}

\textbf{Palabras clave:} cloud lock-in, interoperabilidad, infraestructura como código,
multi-cloud, OpenTofu, FinOps, prueba de concepto, CLI.

\vspace{0.5cm}

\textbf{Resumen:}

La computación en la nube se ha consolidado como el paradigma dominante para el despliegue de
infraestructura, pero su adopción arrastra un problema estructural: el \textit{cloud lock-in},
la dependencia de un proveedor concreto que hace técnica y económicamente costoso migrar a otro.
Este problema se manifiesta en tres dimensiones que se refuerzan mutuamente —la infraestructura,
el conocimiento y el coste—, y tiene su raíz en la falta de interoperabilidad real entre las
plataformas cloud: cada proveedor impone sus propios servicios, formatos y herramientas, de modo
que una misma arquitectura no puede trasladarse de una nube a otra sin reescribirla por completo.

Este Trabajo de Fin de Grado propone \textbf{CloudPorter}, una herramienta de línea de comandos
que aborda el lock-in a nivel de infraestructura. CloudPorter permite describir la infraestructura
una única vez en un manifiesto en formato YAML, independiente del proveedor, y traducirla de forma
determinista a plantillas de \textit{infraestructura como código} con OpenTofu para desplegarla
en distintas nubes. Además de la traducción y el despliegue, la herramienta compara los costes
estimados entre proveedores antes de comprometerse con uno y audita el manifiesto en busca de
anti-patrones de seguridad, todo desde una misma definición agnóstica.

El trabajo se plantea como una \textbf{prueba de concepto} desarrollada de forma iterativa: no
busca entregar un producto terminado, sino construir algo básico pero funcional que demuestre la
viabilidad de la interoperabilidad entre nubes. El resultado es una herramienta con cinco
comandos —validación, auditoría, estimación de costes, traducción y despliegue—, soporte para
recursos de cómputo y de base de datos, compatibilidad con Amazon Web Services y Microsoft Azure,
y una aplicación de demostración de tres capas desplegable en ambas nubes a partir de un único
manifiesto, sin más cambio que el del proveedor de destino.
